\section*{Capítulo \ref{ch:g12:age-of-the-earth} --- La datación radiométrica y la edad de la Tierra}

\begin{solution}{exo:g12:age-of-the-earth:1}
(a) un \omterm{def:g12:age-of-the-earth:closed-system}{sistema cerrado}; (b) una \omterm{def:g12:age-of-the-earth:isochron}{isócrona}; (c) la \omterm{def:g12:age-of-the-earth:age}{edad de la Tierra}.
\end{solution}

\begin{solution}{exo:g12:age-of-the-earth:2}
Plagioclasa (0.35), hornblenda (1.20), roca total (3.10), feldespato K (6.40),
biotita (24.0). La biotita queda más arriba: es la que contiene más rubidio 87
por cada estroncio 86, así que su rubidio ha producido más estroncio 87 desde
que se formó el granito.
\end{solution}

\begin{solution}{exo:g12:age-of-the-earth:3}
$T = \ln 2/\lambda = 0.6931/\num{1.3972e-11} \approx \num{4.96e10}$ años,
49\,600 millones de años: $49.6/4.54 \approx 11$ veces la \omterm{def:g12:age-of-the-earth:age}{edad de la Tierra}.
\end{solution}

\begin{solution}{exo:g12:age-of-the-earth:4}
$\lambda t = \num{1.3972e-11} \times \num{3.10e8} \approx 0.004331$, y
$e^{0.004331} - 1 \approx 0.004341$. Para 480 millones de años, $\lambda t
\approx 0.006707$ y $e^{\lambda t} - 1 \approx 0.006729$. Los dos valores son
cercanos al propio $\lambda t$, porque $\lambda t$ es pequeño.
\end{solution}

\begin{solution}{exo:g12:age-of-the-earth:5}
Phillips (1860), Kelvin (1862), Joly (1899), Holmes (1927), Patterson
(1956). $4540/100 \approx 45$ veces.
\end{solution}

\begin{solution}{exo:g12:age-of-the-earth:6}
$0.7080 + 10.0 \times 0.004341 \approx 0.7514$ hoy; 0.7080 cuando se formó el
granito, como todos los minerales del granito.
\end{solution}

\begin{solution}{exo:g12:age-of-the-earth:7}
(1) $m = (0.77525 - 0.71806)/(9.40 - 0.90) = 0.05719/8.50 \approx 0.006728$;
$t = \ln(1.006728)/\num{1.3972e-11} \approx \num{4.80e8}$ años, 480 millones
de años. (2) $0.71806 + 1.70 \times 0.006728 \approx 0.72950$ y $0.71806 + 4.20
\times 0.006728 \approx 0.74632$: las dos están sobre la recta. (3) Sí: el
granito corta el esquisto, así que es más joven, y 310 es menos que 480
millones de años.
\end{solution}

\begin{solution}{exo:g12:age-of-the-earth:8}
Todos sus puntos caerían en el mismo lugar: un solo punto no define ni una
recta ni una pendiente. La biotita está muy a la derecha de los demás; un
pequeño error en cualquier punto cambia entonces la pendiente de la recta que
pasa por él y por la plagioclasa mucho menos de lo que la cambiaría con dos
puntos muy próximos.
\end{solution}

\begin{solution}{exo:g12:age-of-the-earth:9}
(1) Su rubidio no cambia, así que conserva su $x$; tiene menos estroncio 87,
así que su $y$ baja: el punto se desplaza en vertical hacia abajo. (2) La
recta que pasa por él y por la plagioclasa es menos inclinada, y da una edad
demasiado joven. (3) La biotita ya no está sobre la recta que pasa por los
demás minerales: los puntos no están alineados, señal de que una muestra no ha
sido un \omterm{def:g12:age-of-the-earth:closed-system}{sistema cerrado}.
\end{solution}

\begin{solution}{exo:g12:age-of-the-earth:10}
(1) $4540/96 \approx 47$ veces. (2) Si una parte de la sal sale del mar, por
ejemplo enterrada en sedimentos, el mar contiene menos sal de la que trajeron
en total los ríos: dividir la sal presente por el aporte anual da un tiempo
demasiado corto.
\end{solution}

\begin{solution}{exo:g12:age-of-the-earth:11}
$m = (34.86 - 10.34)/(50.28 - 9.46) = 24.52/40.82 \approx 0.601$, entre las
pendientes de 4550 (0.590) y 4600 millones de años (0.610): unos 4580
millones de años.
\end{solution}

\begin{solution}{exo:g12:age-of-the-earth:12}
(1) El basalto, con la razón más baja, procede del manto; el granito, con la
más alta, al menos en parte de corteza fundida. (2) El rubidio pasa al
fundido cuando el manto se funde, así que se ha concentrado en la corteza, que
ha producido más estroncio 87 con las eras. (3) No: los minerales que
cristalizan a partir de un mismo magma comparten su razón inicial, y estas dos
rocas no.
\end{solution}

\begin{solution}{exo:g12:age-of-the-earth:13}
(1) Rubidio 87: $2^{-4.54/49.61} \approx 0.94$, alrededor del
\qty{94}{\percent}; uranio 238: $2^{-4.54/4.47} \approx 0.49$, más o menos la
mitad. (2) $m =
e^{\lambda t} - 1 \approx \num{1.3972e-11} \times 10^6 \approx
\num{1.4e-5}$: incluso la biotita del granito de Varden ($x = 24$) solo
subiría $0.0003$ por encima de la razón inicial. Se ha formado tan poco
estroncio 87 que la pendiente apenas se distingue de cero: un reloj lento
conviene a las rocas antiguas.
\end{solution}

\begin{solution}{exo:g12:age-of-the-earth:14}
(1) $4.438 + 0.6026 \times 19.0 \approx 15.89$; $15.5$ es menor en $0.39$,
alrededor de un \qty{2.4}{\percent}: el sedimento está cerca de la recta.
(2) 0.6026 está entre 0.590 (4.55) y 0.610 (4.60): unos 4580 millones de
años, dentro del intervalo de Patterson, de 4480 a 4620 millones de años.
\end{solution}

\begin{solution}{exo:g12:age-of-the-earth:15}
(1) La reserva solo pierde entonces $46.7 - 24.4 = \qty{22.3}{TW}$ netos, así
que la misma reserva dura $46.7/22.3 \approx 2.1$ veces más. (2) $96 \times 2.1 \approx
200$ millones de años, todavía más de veinte veces demasiado poco.
(3) La \omterm{def:g12:internal-heat:convection}{convección del manto}, que sube calor desde el interior profundo y
mantiene alto el gradiente cerca de la superficie durante mucho más tiempo que
la conducción sola.
\end{solution}

\begin{solution}{pb:g12:age-of-the-earth:1}
\textbf{1.} La biotita ($x = 24.0$).

\textbf{2.} Cristalizaron a partir de un mismo magma, y los minerales que
cristalizan a partir de un mismo fundido incorporan su estroncio con su razón
de isótopos.

\textbf{3.} El estroncio 86 es estable y no procede de ninguna desintegración:
en un \omterm{def:g12:age-of-the-earth:closed-system}{sistema cerrado} su número no cambia, así que es un patrón fijo.

\textbf{4.} El estroncio 87, producido por la desintegración del rubidio 87;
el rubidio 87 ha disminuido y el estroncio 86 no ha cambiado.

\textbf{5.} Una recta horizontal en $y = 0.7080$: la misma razón, distintos $x$.

\textbf{6.} $(0.81218 - 0.70952)/(24.0 - 0.35) = 0.10266/23.65 \approx
0.004341$.

\textbf{7.} $(0.72146 - 0.70952)/(3.10 - 0.35) = 0.01194/2.75 \approx
0.004342$.

\textbf{8.} Coinciden con un margen de 1 entre 4000. La primera es la más
fiable: sus dos puntos están muy separados, así que el redondeo de las razones
importa menos.

\textbf{9.} $0.70952 - 0.35 \times 0.004341 \approx 0.7080$.

\textbf{10.} $0.7080 + 6.40 \times 0.004341 \approx 0.73578$: exactamente el
valor medido.

\textbf{11.} $e^{\lambda t} = 1 + m = 1.004341$.

\textbf{12.} $\lambda t = \ln 1.004341 \approx 0.004331$.

\textbf{13.} $t = 0.004331/\num{1.3972e-11} \approx \num{3.10e8}$ años.

\textbf{14.} Sí: el granito corta el esquisto, así que, por el principio de
intersección, es más joven, y 310 es menos que 480 millones de años.

\textbf{15.} 310 millones de años está entre 358.9 y 298.9 millones de años:
el Carbonífero.

\textbf{16.} La razón $^{87}\mathrm{Sr}/^{86}\mathrm{Sr}$ del magma cuando
cristalizó el granito.

\textbf{17.} El granito: $0.7080 > 0.7035$.

\textbf{18.} El rubidio pasa al fundido cuando el manto se funde, así que se
ha concentrado en la corteza; con más rubidio, la corteza ha producido más
estroncio 87.

\textbf{19.} Que el magma procedía, al menos en parte, de \omterm{def:g12:continental-crust:continental-crust}{corteza continental}
fundida y no solo del manto.

\textbf{20.} El granito de Varden tiene \textbf{310 millones de años}.
\end{solution}
